\documentclass[a4paper,11pt]{article}

\usepackage[utf8]{inputenc}
\usepackage[margin=0.85in]{geometry}
\usepackage[hidelinks]{hyperref}
\usepackage{parskip}
\usepackage{microtype}

\setlength{\parskip}{5pt}
\setlength{\parindent}{0pt}
\urlstyle{rm}

\begin{document}

%============================ HEADER ============================
\begin{center}
  {\LARGE \scshape \textbf{Md Monirul Islam}}\\[3pt]
  {\large Research Summary}\\[4pt]
  {\small
    \href{mailto:monirulislam.acad@gmail.com}{monirulislam.acad@gmail.com}
    ~$\vert$~ \href{https://Monirul305.github.io}{Portfolio}
    ~$\vert$~ \href{https://scholar.google.com/citations?user=m-6T5FgAAAAJ&hl=en}{Google Scholar}
    ~$\vert$~ \href{https://github.com/Monirul305}{GitHub}
  }
\end{center}
\vspace{2pt}
\hrule
\vspace{8pt}

%============================ INTRO ============================
I work on the reliability of multi-robot systems that operate in dense,
human-occupied environments. Over three years of R\&D at Cybernetics Hi-Tech
Solutions I have designed, built, and field-deployed several such systems --- a
production three-robot AGV fleet, indigenous EOD ground robots, and a mobile
manipulator --- and one theme recurs across all of them: a coordination layer
that is empirically robust but methodologically thin. It works because every
anticipated failure mode was engineered around, not because its safety can be
stated as a guarantee. Closing that gap is what I want to spend a PhD on.
% TAILORING: insert one sentence here naming this professor's work, e.g.
% "Your group's work on <topic> connects directly to the coordination and
% safety questions below."

%======================= CONTRIBUTION 1 ========================
\textbf{Decentralized multi-robot coordination, deployed in production.} My
primary contribution is the coordination scheme behind a three-robot AGV fleet
that has run for months without a collision in a live garment factory, sharing
the floor with hundreds of workers under no physical barriers. It is a
decentralized, prioritized multi-agent path-finding (MAPF) design: each robot
plans its own path with A\textsuperscript{*} and publishes it as timestamped
grid-cell reservations that peers treat as spatiotemporal constraints under
vertex-, swap-, and target-conflict checks, while a fail-closed coordinator
handles dispatch, mutual exclusion, and deadlock recovery. The scheme sits in
the cooperative-A\textsuperscript{*}\,/\,prioritized-planning family, and its
limitations are exactly the questions I find interesting: priorities are
inherited from dispatch order rather than searched, completeness is not
guaranteed, and ``safe'' currently means ``no anticipated failure has
occurred,'' not a characterized bound. (First-author manuscript in preparation.)

%======================= CONTRIBUTION 2 ========================
\textbf{Estimation and the embedded substrate that coordination depends on.}
Reliable coordination rests on layers beneath it that are usually assumed rather
than characterized, and I have built those layers directly. For localization I
designed a five-state extended Kalman filter that fuses wheel odometry with
fiducial camera fixes; velocity-bias states absorb systematic encoder drift over
long, feature-poor corridors, and a rewind-and-replay step applies each delayed
camera measurement at its true timestamp. At the drive-and-comms layer, my
first-author work on GA-tuned bilateral cross-coupled BLDC synchronization drives
the speed difference between two mismatched motors toward zero from a dedicated
inter-MCU exchange, cutting peak inter-motor speed error by 79\,\% and
accumulated heading drift by 80\,\% against an independent-PID baseline (submitted); a co-authored fault-tolerant dual-network feedback scheme keeps a
control loop alive through single-channel failure. These are precisely the
substrate properties a coordination-layer safety claim has to be able to assume.

%======================= CONTRIBUTION 3 ========================
\textbf{Manipulation and further systems.} I also wrote a from-scratch
closed-form inverse-kinematics solver for a six-DOF mobile manipulator ---
replacing a numerical solver with deterministic, branch-stable solutions and
handling the home-pose singularity with a smooth-scaling scheme --- and have
contributed across localization, motion planning, and embedded control on
drones and other unmanned platforms. Two further works are peer-reviewed
(an IEEE ICCIT 2024 paper on multi-node RS-485 coordination and an MDPI
\textit{Micromachines} 2023 paper from my undergraduate device-physics thesis);
the robotics results above are first- or co-author manuscripts currently in
preparation or submitted.

%======================= DIRECTION ========================
\textbf{Research direction.} I want to turn coordination layers that empirically
work into ones with \emph{characterized} safety --- scaled past demonstrative
fleet sizes and characterized down to the embedded substrate they rely on ---
using the formal tools that field-deployed engineering rarely produces on its
own: runtime safety monitors, control-barrier or probabilistic-collision
arguments, and the MAPF completeness and priority-ordering literature. My
strength is building and fielding these systems; the formal apparatus is the
capability I am deliberately seeking a PhD to add. A fuller statement of the
direction and a concrete first-year plan are on my
\href{https://Monirul305.github.io/research-statement/}{research statement} page.
I would welcome the chance to discuss how this fits your group's work.

\end{document}
